\documentclass[11pt]{article}

\usepackage[T1]{fontenc}
\usepackage{lmodern}
\usepackage[margin=1in]{geometry}
\usepackage{amsmath,amssymb,amsthm}
\usepackage{booktabs,longtable}
\usepackage{microtype}
\usepackage[hidelinks]{hyperref}
\setlength{\emergencystretch}{2em}

\newtheorem{theorem}{Theorem}[section]
\newtheorem{lemma}[theorem]{Lemma}

\newcommand{\N}{\mathbb N}

\hypersetup{
  pdftitle={A sparse-defect counterexample to Abelian-border periodicity},
  pdfauthor={Carptopus},
  pdfsubject={A binary counterexample to Abelian-border periodicity},
  pdfkeywords={Abelian border, infinite word, Abelian periodicity, counterexample, combinatorics on words, sparse defects}
}

\title{A Sparse-Defect Counterexample to\\
Abelian-Border Periodicity}
\author{Carptopus\\\small\texttt{carptopus@163.com}}
\date{28 August 2026\\\small Internal candidate; not peer reviewed}

\begin{document}
\maketitle

\begin{abstract}
Charlier, Harju, Puzynina, and Zamboni asked whether an infinite word whose
sufficiently long factors are all Abelian bordered must be Abelian periodic.
We answer the question negatively, already over the binary alphabet. The
construction starts from a four-height periodic nearest-neighbor path whose
every anti-diagonal contains all parity-compatible endpoint sums. Sparse
finite replacement walks then destroy every constant arithmetic-progression
tail while a five-witness redundancy preserves the full local anti-diagonal
palette. The resulting binary word is not Abelian ultimately periodic,
although every factor of length at least $141$ has an Abelian border. The
finite periodic-palette and replacement certificates are exact and
independently reproducible.
\end{abstract}

\noindent\textbf{Keywords.} Abelian border; infinite word; Abelian
periodicity; counterexample; combinatorics on words; sparse defects.

\medskip
\noindent\textbf{2020 Mathematics Subject Classification.} 68R15, 05A05.

\section{Introduction}

A finite word is Abelian bordered if it has a nonempty proper prefix Abelian
equivalent to a suffix. Charlier, Harju, Puzynina, and
Zamboni~\cite[Question~1]{CHPZ} asked whether an infinite word for which every
sufficiently long factor is Abelian bordered must be Abelian periodic. Their
nonperiodic examples retained Abelian periodicity. Fici and Puzynina later
recorded the same question as Problem~42 in their survey~\cite{FP}.

Previous work by the present author determined the first finite-threshold
layers: if every Abelian-unbordered factor has length at most $14$, then the
word is Abelian ultimately periodic with an Abelian period of length at most
$14$, and this bound is sharp~\cite{CarptopusThresholds}. Those results did
not decide what happens under an arbitrary finite bound. The present
construction shows that the unrestricted statement is false.

The mechanism has two parts. First, a periodic height path on
$\{0,1,2,3\}$ supplies every parity-compatible height sum near every center.
Second, widely separated replacement blocks are placed so that every
arithmetic progression is changed infinitely often. Five translated
witnesses for each center and sum ensure that a single replacement block
cannot destroy them all.

The result is not a statement about the minimum possible threshold, period,
or height range. The constants $28$ and $141$ are explicit parameters of one
construction. Recent work on mutually Abelian-bordered finite
words~\cite{MaityKrishna} and on Abelian periodicity of purely morphic
words~\cite{FilimonovaPuzynina} concerns different object classes.

\section{Definitions and the height-path translation}

Two finite words are \emph{Abelian equivalent} if they have the same Parikh
vector; in particular, Abelian-equivalent words have the same length. For
binary words of the same length, this is equivalent to containing the same
number of $1$s. A finite word is \emph{Abelian bordered} if it has a nonempty
proper prefix Abelian equivalent to a suffix. An infinite word $w$ is
\emph{Abelian ultimately periodic} if
\[
w=u v_0v_1v_2\cdots,
\]
where the nonempty finite words $v_n$ are pairwise Abelian equivalent.
Following~\cite{CHPZ,FP}, we use \emph{Abelian periodic} with this allowed
finite preperiod.

For a binary word $w=w_0w_1\cdots$, define
\[
z_0=0,\qquad z_{n+1}-z_n=2w_n-1. \tag{2.1}
\]
Then $z$ is an integer nearest-neighbor path and
\[
|w[u,v)|_1=\frac{v-u+z_v-z_u}{2}. \tag{2.2}
\]

\begin{lemma}[Center contraction]\label{lem:center}
Fix $C\ge1$. Every factor of $w$ of length at least $C$ is Abelian bordered
if and only if, for all $0\le i<j$, there exist $i\le a\le b\le j$ such that
\[
a+b=i+j,\qquad b-a<C,\qquad z_a+z_b=z_i+z_j. \tag{2.3}
\]
\end{lemma}

\begin{proof}
If $j-i<C$, take $(a,b)=(i,j)$. Suppose $j-i\ge C$. An Abelian border can
be chosen with length $1\le k\le(j-i)/2$: if a longer prefix and suffix
overlap, deleting their common overlap gives the complementary shorter
border. By (2.2), equality of the two Parikh vectors is equivalent to
\[
z_{i+k}+z_{j-k}=z_i+z_j.
\]
Replace $(i,j)$ by $(i+k,j-k)$ and repeat while the new gap is at least $C$.
The gap strictly decreases, whereas the center and endpoint sum remain fixed,
so (2.3) follows.

Conversely, if $j-i\ge C$, a pair supplied by (2.3) cannot be the original
endpoint pair. Writing $k=a-i=j-b>0$ and applying (2.2) gives an Abelian
border of length $k$.
\end{proof}

\begin{lemma}[Arithmetic-progression criterion]\label{lem:ap}
Assume that $z$ is bounded. The word $w$ is Abelian ultimately periodic if
and only if there exist $A\ge0$ and $h\ge1$ such that
\[
z_{A+nh}\quad(n\ge0)
\]
is constant.
\end{lemma}

\begin{proof}
If the length-$h$ blocks beginning at $A+nh$ have one common Parikh vector,
(2.2) says that
\[
z_{A+(n+1)h}-z_{A+nh}=\delta
\]
is independent of $n$. Boundedness forces $\delta=0$. Conversely, a constant
arithmetic-progression tail makes all those length-$h$ blocks have the same
number of $1$s by (2.2).
\end{proof}

\section{A periodic full-palette skeleton}

Set $d=28$, and extend the height word
\[
P=0101010101232323210123232321 \tag{3.1}
\]
periodically to $p\colon\N_0\to\{0,1,2,3\}$. Consecutive heights, including
the last and first entries of $P$, differ by $1$.

For a center residue $m\in\mathbb Z/d\mathbb Z$, define
\[
\Pi_m=\{p_r+p_{m-r}:r\in\mathbb Z/d\mathbb Z\}. \tag{3.2}
\]

\begin{lemma}[Full palette]\label{lem:palette}
For every $m\in\mathbb Z/d\mathbb Z$,
\[
\Pi_m=
\begin{cases}
\{0,2,4,6\},&m\equiv0\pmod2,\\
\{1,3,5\},&m\equiv1\pmod2.
\end{cases} \tag{3.3}
\]
\end{lemma}

\begin{proof}
Substitution of the $28$ entries of (3.1) into (3.2) gives (3.3). The exact
verification program checks equality, rather than mere containment, for all
$28$ center residues.
\end{proof}

We also need to change the skeleton at an arbitrary residue without breaking
the nearest-neighbor path at the boundary of one period block.

\begin{lemma}[Replacement walks]\label{lem:replacement}
For every target residue $t\in\mathbb Z/d\mathbb Z$, there exist a start
residue $s_t$, an offset $1\le r_t<d$, and a word
$Q^{(t)}=q_0q_1\cdots q_{d-1}$ over $\{0,1,2,3\}$ such that
\begin{enumerate}
\item $t\equiv s_t+r_t\pmod d$;
\item $q_0=p_{s_t}$ and $|q_{d-1}-q_0|=1$;
\item $|q_{j+1}-q_j|=1$ for $0\le j<d-1$;
\item $q_{r_t}\ne p_t$.
\end{enumerate}
\end{lemma}

\begin{proof}
The $28$ explicit certificates are listed in Appendix~\ref{app:certificates}.
The supplementary verifier checks them directly; its deterministic recursive
depth-first search also regenerates one valid certificate for every target
residue. Conditions 2 and 3 ensure that replacing one aligned period block by
$Q^{(t)}$ preserves the nearest-neighbor condition both inside the block and
across its two boundaries.
\end{proof}

\section{Sparse defects}

Enumerate the pairs $(a,h)\in\N_0\times\N$ in a sequence
$(a_j,h_j)_{j\ge1}$ in which every pair occurs infinitely many times. At
stage $j$, choose a target
\[
t_j\in\{a_j+n h_j:n\ge0\}
\]
far enough to the right that the aligned replacement interval $I_j$ supplied
by Lemma~\ref{lem:replacement} lies after all previous intervals, the distance
between successive intervals is greater than $6d$, and these distances tend
to infinity. This is possible because every arithmetic progression is
unbounded and the target offset inside a replacement block is less than $d$.

On $I_j$, replace $p$ by $Q^{(t_j\bmod d)}$; outside all replacement
intervals retain $p$. Lemma~\ref{lem:replacement} gives a nearest-neighbor
path
\[
z\colon\N_0\longrightarrow\{0,1,2,3\}
\]
such that $z_{t_j}\ne p_{t_j}$ for every $j$. Let $D$ be the set of
positions where $z\ne p$.

\begin{lemma}\label{lem:density}
The set $D$ has natural density zero, and every arithmetic progression
$\{a+nh:n\ge0\}$ contains infinitely many designated targets $t_j$.
\end{lemma}

\begin{proof}
The second assertion follows from the repeated enumeration. The intervals
have one fixed length and their successive gaps tend to infinity. Given
$G>0$, all but finitely many consecutive intervals are separated by at least
$G$; hence their upper density is at most $d/(d+G)$. Letting $G\to\infty$
proves that the interval union, and therefore $D$, has density zero.
\end{proof}

\section{Preservation of center contraction}

\begin{lemma}\label{lem:preservation}
The path $z$ satisfies (2.3) with $C=141$.
\end{lemma}

\begin{proof}
Let $0\le i<j$ with $j-i\ge141$, set $N=i+j$, and put $v=z_i+z_j$.
Because every nearest-neighbor path has fixed height parity, $v$ is one of
the values allowed by the corresponding case of (3.3). Choose $r\pmod d$
such that
\[
p_r+p_{N-r}=v.
\]
Among the pairs $(a,b)$ satisfying $a+b=N$ and $a\equiv r\pmod d$, choose
$(a_0,b_0)$ with $|a_0-b_0|\le d$. Consider the five pairs
\[
(a_k,b_k)=(a_0+kd,b_0-kd),\qquad -2\le k\le2. \tag{5.1}
\]
They have the same center and the same skeleton sum $v$, and
\[
|a_k-b_k|\le5d=140. \tag{5.2}
\]
Since $N\ge141$, all endpoints in (5.1) are nonnegative; exchanging the two
entries of a pair if necessary does not change its properties.

The candidate endpoints occupy a band of span at most $5d$. Since distinct
replacement intervals are separated by more than $6d$, this band meets at
most one of them. A replacement interval of length $d$ meets at most two
members of the progression $(a_k)$ and at most two members of $(b_k)$. It
therefore contaminates at most four of the five pairs. For a remaining clean
pair,
\[
z_{a_k}+z_{b_k}=p_{a_k}+p_{b_k}=v=z_i+z_j.
\]
Together with (5.2), this is (2.3). Because the pair has center $N/2$, its
gap is smaller than the original gap, so it lies inside $[i,j]$.
\end{proof}

\section{Destruction of every Abelian period}

\begin{lemma}\label{lem:destroy}
The path $z$ is not constant on the tail of any arithmetic progression.
\end{lemma}

\begin{proof}
Fix $A\ge0$ and $h\ge1$. The periodic sequence $p_{A+nh}$ is either
constant or takes at least two values periodically.

In the second case, each of two distinct values occurs on a subprogression of
positive natural density in $\N_0$. A zero-density set cannot contain a tail
of either subprogression. Hence $z_{A+nh}$ takes both values infinitely often
outside $D$.

In the first case, the repeated enumeration designates infinitely many
distinct targets on the progression, and each target is changed from its
skeleton value. The progression also contains infinitely many points outside
$D$, because it has positive density while $D$ has density zero. Thus
$z_{A+nh}$ is again not eventually constant.
\end{proof}

\section{Main theorem}

\begin{theorem}\label{thm:main}
There exists a binary infinite word that is not Abelian ultimately periodic
but whose every factor of length at least $141$ is Abelian bordered.
Consequently, the answer to Question~1 of~\cite{CHPZ} and Problem~42
of~\cite{FP} is negative.
\end{theorem}

\begin{proof}
Define
\[
w_n=\frac{1+z_{n+1}-z_n}{2}.
\]
The nearest-neighbor property makes $w_n\in\{0,1\}$. Lemmas
\ref{lem:preservation} and~\ref{lem:center} show that every factor of length
at least $141$ is Abelian bordered. Lemmas~\ref{lem:destroy} and~\ref{lem:ap}
show that $w$ is not Abelian ultimately periodic.
\end{proof}

\section{Remarks}

\begin{enumerate}
\item Theorem~\ref{thm:main} does not assert that $141$ is the least possible
threshold.
\item The period $28$ and the four-height range are not claimed to be
minimal.
\item The finite verification discovers and checks the skeleton certificates,
but the infinite sparse-set and arithmetic-progression arguments are
mathematical proofs independent of finite-prefix testing.
\item No morphic representation of the construction is asserted.
Reference~\cite{FilimonovaPuzynina} treats the restricted class of purely
morphic words and does not cover the unrestricted hypothesis of
Theorem~\ref{thm:main}.
\end{enumerate}

\appendix
\section{Replacement certificates}\label{app:certificates}

Each row gives $(t,s_t,r_t,Q^{(t)})$. The word has length $28$.

\begingroup
\scriptsize
\setlength{\tabcolsep}{5pt}
\begin{longtable}{rrr l}
\toprule
$t$ & $s_t$ & $r_t$ & $Q^{(t)}$\\
\midrule
\endfirsthead
\toprule
$t$ & $s_t$ & $r_t$ & $Q^{(t)}$\\
\midrule
\endhead
0 & 1 & 27 & \texttt{1010101010101010101010101012}\\
1 & 3 & 26 & \texttt{1010101010101010101010101232}\\
2 & 0 & 2 & \texttt{0121010101010101010101010101}\\
3 & 0 & 3 & \texttt{0123210101010101010101010101}\\
4 & 0 & 4 & \texttt{0101210101010101010101010101}\\
5 & 0 & 5 & \texttt{0101232101010101010101010101}\\
6 & 0 & 6 & \texttt{0101012101010101010101010101}\\
7 & 0 & 7 & \texttt{0101012321010101010101010101}\\
8 & 0 & 8 & \texttt{0101010121010101010101010101}\\
9 & 0 & 9 & \texttt{0101010123210101010101010101}\\
10 & 0 & 10 & \texttt{0101010101010101010101010101}\\
11 & 0 & 11 & \texttt{0101010101010101010101010101}\\
12 & 0 & 12 & \texttt{0101010101010101010101010101}\\
13 & 0 & 13 & \texttt{0101010101010101010101010101}\\
14 & 0 & 14 & \texttt{0101010101010101010101010101}\\
15 & 0 & 15 & \texttt{0101010101010101010101010101}\\
16 & 0 & 16 & \texttt{0101010101010101010101010101}\\
17 & 0 & 17 & \texttt{0101010101010101232101010101}\\
18 & 0 & 18 & \texttt{0101010101010101012101010101}\\
19 & 0 & 19 & \texttt{0101010101010101012321010101}\\
20 & 0 & 20 & \texttt{0101010101010101010101010101}\\
21 & 0 & 21 & \texttt{0101010101010101010101010101}\\
22 & 0 & 22 & \texttt{0101010101010101010101010101}\\
23 & 0 & 23 & \texttt{0101010101010101010101010101}\\
24 & 0 & 24 & \texttt{0101010101010101010101010101}\\
25 & 0 & 25 & \texttt{0101010101010101010101010101}\\
26 & 0 & 26 & \texttt{0101010101010101010101010101}\\
27 & 1 & 26 & \texttt{1010101010101010101010101232}\\
\bottomrule
\end{longtable}
\endgroup

\section{Reproducibility}

The supplementary file \path{verify_sparse_palette_counterexample.py} has
SHA-256
\begin{center}
\small\ttfamily
92628A23B7ACED4E111A9909C0EBC429\\
07981737479AB902240CABC0B61B7ECC.
\end{center}
The program checks the cyclic nearest-neighbor skeleton, all $28$ palette equalities, all
$28$ replacement certificates, the four-of-five contamination bound, and a
finite sparse-prefix positive control. Its role and limitations are stated in
the program header.

\section*{AI-assisted research and writing disclosure}

The author used OpenAI Codex for computer-assisted exploration, exact
verification, proof checking, literature triage, and drafting support. The
author reviewed the claims and assumes responsibility for the content.

\begin{thebibliography}{9}

\bibitem{CHPZ}
\'E. Charlier, T. Harju, S. Puzynina, and L. Q. Zamboni,
Abelian bordered factors and periodicity,
\emph{European Journal of Combinatorics} 51 (2016), 407--418.
\href{https://doi.org/10.1016/j.ejc.2015.07.003}{doi:10.1016/j.ejc.2015.07.003}.

\bibitem{FP}
G. Fici and S. Puzynina,
Abelian combinatorics on words: a survey,
\emph{Computer Science Review} 47 (2023), 100532.
\href{https://doi.org/10.1016/j.cosrev.2022.100532}{doi:10.1016/j.cosrev.2022.100532}.

\bibitem{CarptopusThresholds}
Carptopus,
Sharp Initial Thresholds for Abelian-Bordered Binary Infinite Words,
Public Beta v0.1-beta (2026).
\href{https://doi.org/10.5281/zenodo.22100526}{doi:10.5281/zenodo.22100526}.

\bibitem{MaityKrishna}
A. Maity and K. V. Krishna,
Mutually Abelian-Bordered Binary Words,
\href{https://arxiv.org/abs/2509.20773}{arXiv:2509.20773v1} (2025).

\bibitem{FilimonovaPuzynina}
A. Filimonova and S. Puzynina,
On abelian periodicity of purely morphic words,
\href{https://arxiv.org/abs/2605.30306}{arXiv:2605.30306v1} (2026).

\end{thebibliography}

\end{document}
